\documentclass[11pt,a4paper]{article}

\usepackage[utf8]{inputenc}
\usepackage[margin=0.9in]{geometry}
\usepackage{amsmath,amssymb,amsthm,amsfonts}
\usepackage{mathtools}
\usepackage{hyperref}
\usepackage{booktabs}
\usepackage{cite}
\usepackage{microtype}
\usepackage{array}

\hypersetup{
    colorlinks=true,
    linkcolor=blue,
    citecolor=red,
    urlcolor=blue
}

% Theorem Environments
\theoremstyle{plain}
\newtheorem{theorem}{Theorem}[section]
\newtheorem{lemma}[theorem]{Lemma}
\newtheorem{proposition}[theorem]{Proposition}
\newtheorem{corollary}[theorem]{Corollary}

\theoremstyle{definition}
\newtheorem{definition}[theorem]{Definition}
\newtheorem{example}[theorem]{Example}

\theoremstyle{remark}
\newtheorem{remark}[theorem]{Remark}

\numberwithin{equation}{section}

\title{\textbf{Inter-Dimensional Simplicial Transforms, Fractional Boundary Traces, and Matrix Beta-Kernels}}
\author{\textbf{Reinaldo Maia Silva-Filho} \\ \textit{Universidade Federal de Lavras (UFLA)}}
\date{\today}


\DeclareMathOperator*{\argmin}{argmin}
\providecommand{\Hyp}{\mathbb{H}}
\providecommand{\Sph}{\mathbb{S}}
\providecommand{\II}{\mathrm{II}}
\providecommand{\R}{\mathbb{R}}
\providecommand{\Area}{\mathrm{Area}}
\providecommand{\Hsn}{\mathcal{H}}
\providecommand{\BV}{\mathrm{BV}}
\providecommand{\TV}{\mathrm{TV}}
\providecommand{\cutnorm}[1]{\|#1\|_{\square}}

\begin{document}

\maketitle

\begin{abstract}
We develop a comprehensive mathematical theory of dimension-changing non-local operators $\mathcal{T}_{m \to n}^\alpha : L^p(\mathbb{R}^m) \to L^q(\mathbb{R}^n)$ induced by continuous multinomial Beta-kernels on simplices. For dimensional reduction ($m > n$), we introduce the \emph{Fractional Radon--Beta Transform} $\mathcal{R}_{m \to n}^\alpha$ and prove the Sobolev trace estimate $\mathcal{R}_{m \to n}^\alpha : H^s(\mathbb{R}^m) \to H^{s + \gamma - (m-n)/2}(\mathbb{R}^n)$, where $\gamma$ is the Fourier decay rate of the kernel. We show that for the Beta-kernel $\gamma = 1$ exactly, independently of $\alpha$, and that no choice of order makes the trace an isomorphism. For dimensional elevation ($m < n$), the dual \emph{Simplicial Extension Operator} $\mathcal{E}_{m \to n}^\alpha$ loses $(n-m)/2$ derivatives, up to the kernel gain.

We discuss four applications: (i) \textbf{Coupled Multiscale PDEs:} a coupled 3D bulk / 2D surface / 1D filament transport system with mass conservation and energy dissipation, as a formal computation under an explicit geometric condition on the interfaces; (ii) \textbf{Reconstruction:} the transform restricts a smoothed field to planes through the origin, so the data determine the smoothed field directly, and recovering $f$ reduces to a deconvolution, which is ill-posed and requires regularization; (iii) \textbf{Barycentric Maps:} a Lipschitz bound, with sharp constant $\alpha/(\alpha+m)$, for an affine map of barycentric coordinates defined by Dirichlet posterior means; no lower (bi-Lipschitz) bound holds; and (iv) \textbf{Matrix Beta Operators:} we recall the Siegel--Wishart matrix Beta operators on the cone $\operatorname{Sym}_m^+(\mathbb{R})$ and their zonal-polynomial eigenfunctions.
\end{abstract}

\tableofcontents
\newpage

% ==============================================================================
\section{Introduction and Problem Formulation}
% ==============================================================================

In modern analysis and mathematical physics, linear differential and integral operators predominantly map within a fixed-dimensional space ($L^p(\mathbb{R}^d) \to L^q(\mathbb{R}^d)$). However, many cutting-edge physical systems inherently couple phenomena across distinct spatial dimensions:
\begin{enumerate}
    \item \textbf{Multiscale Interface Dynamics:} Mass, heat, and charge exchange between a 3D bulk fluid, a 2D elastic boundary membrane, and 1D embedded filamentary networks (e.g., graphene on substrates, cellular electrophysiology, and fluid-structure interactions).
    \item \textbf{Integral Geometry and Medical Imaging:} The reconstruction of higher-dimensional $m$-dimensional volume fields from lower-dimensional $n$-dimensional planar projection scans ($n < m$).
    \item \textbf{Geometric Data Science on Hypergraphs:} Dimensionality reduction of higher-order multi-way relational structures (simplicial complexes) into Euclidean visual spaces.
\end{enumerate}

Classical fractional calculus operates via endomorphic Riesz multipliers $|\mathbf{k}|^\alpha$ on $\mathbb{R}^d$ \cite{samko1993}. In this paper, we construct and analyze an inter-dimensional family of non-local operators based on the continuous multinomial Beta-kernel on the standard $(m-1)$-simplex \cite{silvafilho2026simplex, silvafilho2026waves}:
\begin{equation}
\Delta_{m-1}(\alpha) \coloneqq \left\{\mathbf{y} \in \mathbb{R}_{\ge 0}^{m-1} \;\middle|\; \sum_{j=1}^{m-1} y_j \le \alpha\right\}, \quad \binom{\alpha}{\mathbf{y}} = \frac{\Gamma(\alpha+1)}{\prod_{j=1}^m \Gamma(y_j+1)}, \quad y_m \coloneqq \alpha - \sum_{j=1}^{m-1} y_j .
\end{equation}
In Section~2 the kernel is used one dimension higher: on $\Delta_m(\alpha) \subset \mathbb{R}^m$, with $m+1$ parts $y_1, \dots, y_m$ and $y_{m+1} = \alpha - \sum_{j \le m} y_j$, so that it lives in the same space $\mathbb{R}^m$ as the functions it acts on.

% ==============================================================================
\section{The Fractional Radon--Beta Transform \texorpdfstring{($\mathbb{R}^m \to \mathbb{R}^n$)}{(Rm -> Rn)}}
% ==============================================================================

Let $m > n \ge 1$ be positive integers. Let $\mathbf{P} \in \mathbb{R}^{n \times m}$ be a full-rank surjective linear projection matrix ($\operatorname{rank}(\mathbf{P}) = n$). Its Moore--Penrose pseudoinverse $\mathbf{P}^+ \in \mathbb{R}^{m \times n}$ is given by
\begin{equation}
\mathbf{P}^+ \coloneqq \mathbf{P}^T (\mathbf{P} \mathbf{P}^T)^{-1}.
\end{equation}

\begin{definition}[Fractional Radon--Beta Transform]
\label{def:radon_beta}
Generalizing classical integral geometry and Radon projections \cite{helgason2011}, let $\mathcal{K}_\alpha$ denote the normalized continuous multinomial Beta-kernel of order $\alpha > 0$ on the $m$-dimensional simplex $\Delta_m(\alpha) \subset \mathbb{R}^m_{\ge 0}$, so that the kernel lives in the same space as $f$. For $f \in \mathcal{S}(\mathbb{R}^m)$, the Fractional Radon--Beta Transform $\mathcal{R}_{m \to n}^\alpha : \mathcal{S}(\mathbb{R}^m) \to C^\infty(\mathbb{R}^n)$ is
\begin{equation}
\label{eq:radon_beta_def}
\mathcal{R}_{m \to n}^\alpha f(\mathbf{z}) \coloneqq \int_{\Delta_m(\alpha)} \mathcal{K}_\alpha(\mathbf{y})\, f\left(\mathbf{P}^+ \mathbf{z} + \mathbf{y}\right) d\mathbf{y}, \quad \mathbf{z} \in \mathbb{R}^n,
\end{equation}
that is, convolution with the reflected kernel followed by restriction to the $n$-plane $\operatorname{ran}(\mathbf{P}^+)$. In particular $\mathcal{R}_{m\to n}^\alpha(1) \equiv 1$.
\end{definition}

\begin{proposition}[Fourier Representation]
\label{prop:fourier_rep}
Let $\boldsymbol{\xi} \in \mathbb{R}^n$ denote the dual frequency coordinates. Then
\begin{equation}
\widehat{\mathcal{R}_{m \to n}^\alpha f}(\boldsymbol{\xi}) = (2\pi)^{n-m} \sqrt{\det(\mathbf{P}\mathbf{P}^T)} \int_{\operatorname{ker}(\mathbf{P})} \widehat{f}\left(\mathbf{P}^T \boldsymbol{\xi} + \boldsymbol{\eta}\right) \widehat{\mathcal{K}}_\alpha\left(-\mathbf{P}^T \boldsymbol{\xi} - \boldsymbol{\eta}\right) d\boldsymbol{\eta},
\end{equation}
where $\widehat{\mathcal{K}}_\alpha(\mathbf{k}) = \int \mathcal{K}_\alpha(\mathbf{y}) e^{-i\mathbf{k}\cdot\mathbf{y}}d\mathbf{y}$ is the Fourier transform of the kernel, with $|\widehat{\mathcal{K}}_\alpha| \le 1$. (It is not given by the lattice closed form $(\frac{1}{m}\sum_j e^{-ik_j})^\alpha$; see Chapter~3, Remark~6.5.)
\end{proposition}

\begin{proof}
Applying the definition of the Fourier transform on $\mathbb{R}^n$ and expanding $f$ via the inverse Fourier transform on $\mathbb{R}^m$:
\begin{equation}
\widehat{\mathcal{R}_{m \to n}^\alpha f}(\boldsymbol{\xi}) = \frac{1}{(2\pi)^m}\int_{\mathbb{R}^m} \widehat{f}(\mathbf{k}) \left[\int_{\mathbb{R}^n} e^{i((\mathbf{P}^+)^T\mathbf{k} - \boldsymbol{\xi})\cdot\mathbf{z}} d\mathbf{z}\right] \left[\int \mathcal{K}_\alpha(\mathbf{y}) e^{i\mathbf{k}\cdot\mathbf{y}} d\mathbf{y}\right] d\mathbf{k}.
\end{equation}
The $\mathbf{z}$-integral yields $(2\pi)^n \delta((\mathbf{P}^+)^T\mathbf{k} - \boldsymbol{\xi})$, enforcing $\mathbf{k} = \mathbf{P}^T\boldsymbol{\xi} + \boldsymbol{\eta}$ for $\boldsymbol{\eta} \in \operatorname{ker}(\mathbf{P})$, with Jacobian factor $(2\pi)^{n-m}\sqrt{\det(\mathbf{P}\mathbf{P}^T)}$. The $\mathbf{y}$-integral is $\widehat{\mathcal{K}}_\alpha(-\mathbf{k})$.
\end{proof}

% ==============================================================================
\section{Fractional Sobolev Trace and Extension Theorems}
% ==============================================================================

\begin{theorem}[Sobolev Trace Theorem with Kernel Smoothing]
\label{thm:sobolev_trace}
Let $m > n$ and suppose the kernel satisfies the Fourier decay bound
\begin{equation}
\label{eq:decay_gamma}
|\widehat{\mathcal{K}}_\alpha(\mathbf{k})| \le C (1 + |\mathbf{k}|^2)^{-\gamma/2}, \qquad \mathbf{k} \in \mathbb{R}^m,
\end{equation}
for some $\gamma \ge 0$. Then for every $s \in \mathbb{R}$ with $s + \gamma > \frac{m-n}{2}$,
\begin{equation}
\mathcal{R}_{m \to n}^\alpha : H^s(\mathbb{R}^m) \longrightarrow H^{s + \gamma - \frac{m-n}{2}}(\mathbb{R}^n)
\end{equation}
is bounded. Since $|\widehat{\mathcal{K}}_\alpha| \le 1$, the case $\gamma = 0$ always applies and reproduces the classical trace loss of $\frac{m-n}{2}$ derivatives \cite{adams2003}.
\end{theorem}

\begin{proof}
By the Cauchy--Schwarz inequality on the fiber $\boldsymbol{\eta} \in \operatorname{ker}(\mathbf{P}) \cong \mathbb{R}^{m-n}$, together with \eqref{eq:decay_gamma},
\[
\left|\widehat{\mathcal{R}_{m \to n}^\alpha f}(\boldsymbol{\xi})\right|^2 \le C \left(\int_{\mathbb{R}^{m-n}} (1 + |\mathbf{P}^T\boldsymbol{\xi} + \boldsymbol{\eta}|^2)^{-s - \gamma} d\boldsymbol{\eta}\right) \int_{\mathbb{R}^{m-n}} (1 + |\mathbf{P}^T\boldsymbol{\xi} + \boldsymbol{\eta}|^2)^s |\widehat{f}(\mathbf{P}^T\boldsymbol{\xi} + \boldsymbol{\eta})|^2 d\boldsymbol{\eta}.
\]
Since $\mathbf{P}^T\boldsymbol{\xi} \perp \ker\mathbf{P}$ and $|\mathbf{P}^T\boldsymbol{\xi}| \ge \sigma_{\min}(\mathbf{P})|\boldsymbol{\xi}|$, the first factor is at most $C(1 + |\boldsymbol{\xi}|^2)^{-(s+\gamma) + \frac{m-n}{2}}$ when $s + \gamma > \frac{m-n}{2}$. Multiplying by $(1+|\boldsymbol{\xi}|^2)^{s + \gamma - \frac{m-n}{2}}$ and integrating over $\boldsymbol{\xi}$ reassembles $\|f\|_{H^s(\mathbb{R}^m)}^2$ by the orthogonal decomposition $\mathbb{R}^m = \operatorname{ran}\mathbf{P}^T \oplus \ker\mathbf{P}$.
\end{proof}

\begin{remark}[What decay rate $\gamma$ does the Beta-kernel have?]
\label{rem:gamma}
In general $\gamma \neq \alpha$: the smoothing gain is not the fractional order. Nor is there an ``isomorphism'' $H^s(\mathbb{R}^m) \to H^s(\mathbb{R}^n)$ at $\alpha = \frac{m-n}{2}$, since traces are never injective.
(a) The continuous Beta-kernel is strictly positive up to the facets of the simplex. It is a smooth function times the indicator of a polytope, so its Fourier transform decays no faster than $|\mathbf{k}|^{-1}$ in facet-normal directions. Hence $\gamma \le 1$ regardless of $\alpha$.
(b) A trace operator onto a lower-dimensional plane is never injective, because every function vanishing near the plane has zero trace. So no choice of order gives an isomorphism. What $\gamma = \frac{m-n}{2}$ would give is an \emph{index-preserving bounded} map; this requires $\frac{m-n}{2} \le 1$ for the Beta-kernel.
(c) The bound \eqref{eq:decay_gamma} holds with $\gamma = 1$, so $\gamma = 1$ is the exact isotropic exponent. Write $\mathcal{K}_\alpha = \kappa\,\mathbf{1}_{\Delta_m(\alpha)}$, where $\kappa(\mathbf{y}) = \Gamma(\alpha+1)/(I\,\prod_{j=1}^{m+1}\Gamma(y_j+1))$ is $C^\infty$ on $\mathbb{R}^m$ and $I$ is the normalization. Given $\mathbf{k} \ne \mathbf{0}$, choose $j$ with $|k_j| \ge |\mathbf{k}|/\sqrt{m}$ and integrate by parts in $y_j$ using $e^{-i\mathbf{k}\cdot\mathbf{y}} = \frac{i}{k_j}\partial_{y_j}e^{-i\mathbf{k}\cdot\mathbf{y}}$ (divergence theorem on the polytope):
\[
|\widehat{\mathcal{K}}_\alpha(\mathbf{k})| \le \frac{1}{|k_j|}\Big( \|\kappa\|_{L^\infty(\partial\Delta_m(\alpha))}\, \mathcal{H}^{m-1}(\partial\Delta_m(\alpha)) + \|\partial_{y_j}\kappa\|_{L^1(\Delta_m(\alpha))} \Big) \le \frac{C_\alpha}{|\mathbf{k}|}.
\]
Together with $|\widehat{\mathcal{K}}_\alpha| \le 1$ this gives \eqref{eq:decay_gamma} with $\gamma = 1$. By (a), no larger $\gamma$ is possible. Consequently $\mathcal{R}^\alpha_{m \to n} \colon H^s(\mathbb{R}^m) \to H^{s + 1 - \frac{m-n}{2}}(\mathbb{R}^n)$ for $s > \frac{m-n}{2} - 1$. In particular $\mathcal{R}^\alpha_{3 \to 2}$ is bounded from $L^2(\mathbb{R}^3)$ to $H^{1/2}(\mathbb{R}^2) \subset L^2(\mathbb{R}^2)$.
\end{remark}

\begin{theorem}[Dual Simplicial Extension Operator]
For dimensional elevation ($m < n$), let $\mathcal{E}_{m \to n}^\alpha \coloneqq (\mathcal{R}_{n \to m}^\alpha)^*$ be the $L^2$-adjoint. Under \eqref{eq:decay_gamma} for the kernel on $\mathbb{R}^n$, for every $t < 0$ it is bounded as
\begin{equation}
\mathcal{E}_{m \to n}^\alpha : H^{t}(\mathbb{R}^m) \longrightarrow H^{t + \gamma - \frac{n-m}{2}}(\mathbb{R}^n).
\end{equation}
\end{theorem}
\begin{proof}
By Theorem~\ref{thm:sobolev_trace} with the roles of the dimensions exchanged,
\[
\mathcal{R}_{n\to m}^\alpha : H^{s}(\mathbb{R}^n) \to H^{s+\gamma-\frac{n-m}{2}}(\mathbb{R}^m) \qquad \text{for } s + \gamma > \tfrac{n-m}{2}.
\]
Its adjoint maps $H^{-(s+\gamma-\frac{n-m}{2})}(\mathbb{R}^m) \to H^{-s}(\mathbb{R}^n)$. Setting $t = -(s + \gamma - \frac{n-m}{2}) < 0$ gives the claim. Extension thus \emph{loses} $\frac{n-m}{2}$ derivatives up to the kernel gain $\gamma$, as expected: the adjoint of a trace spreads a function on an $m$-plane as a singular layer in $\mathbb{R}^n$.
\end{proof}

\section{Coupled Multiscale 3D--2D--1D Interface Dynamical Systems}
% ==============================================================================

Consider a multiscale coupled physical system consisting of a 3D bulk domain $\Omega \subset \mathbb{R}^3$, an active 2D interface membrane $\Gamma \subset \mathbb{R}^2$, and a 1D conductive filament $\Sigma \subset \mathbb{R}^1$, subject to homogeneous Neumann boundary conditions. Let $u(\mathbf{x}, t)$, $v(\mathbf{y}, t)$, and $w(z, t)$ denote the respective state fields.

\begin{definition}[Coupled Inter-Dimensional Transport System]
Adapting transmission boundary conditions and multiscale interface parabolic systems \cite{evans2010}, we formulate the transmission-coupled dynamical system:
\begin{equation}
\label{eq:coupled_system}
\begin{cases}
\partial_t u(\mathbf{x}, t) = \Delta u(\mathbf{x}, t) - \mathcal{E}_{2 \to 3}^\alpha\big(\lambda (\mathcal{R}_{3 \to 2}^\alpha u - v)\big)(\mathbf{x}, t), & \mathbf{x} \in \Omega \subset \mathbb{R}^3, \\
\partial_t v(\mathbf{y}, t) = \Delta_{\Gamma} v(\mathbf{y}, t) + \lambda \big(\mathcal{R}_{3 \to 2}^\alpha u - v\big)(\mathbf{y}, t) - \mathcal{E}_{1 \to 2}^\beta\big(\kappa (\mathcal{R}_{2 \to 1}^\beta v - w)\big)(\mathbf{y}, t), & \mathbf{y} \in \Gamma \subset \mathbb{R}^2, \\
\partial_t w(z, t) = \partial_z^2 w(z, t) + \kappa \big(\mathcal{R}_{2 \to 1}^\beta v - w\big)(z, t), & z \in \Sigma \subset \mathbb{R}^1,
\end{cases}
\end{equation}
where $\lambda, \kappa > 0$ are transmission coefficients and $\alpha, \beta > 0$ are fixed kernel orders. Here $\Gamma$ is identified with a subset of the plane $\operatorname{ran}(\mathbf{P}^+) \subset \mathbb{R}^3$ used by $\mathcal{R}_{3\to 2}^\alpha$, and $\Sigma$ with a subset of the corresponding line in $\mathbb{R}^2$. We assume the interface condition
\begin{equation}
\label{eq:interface_condition}
\mathbf{P}^+\mathbf{y} + \Delta_3(\alpha) \subset \Omega \quad (\mathbf{y} \in \Gamma), \qquad \mathbf{Q}^+ z + \Delta_2(\beta) \subset \Gamma \quad (z \in \Sigma),
\end{equation}
where $\mathbf{Q}$ is the projection used by $\mathcal{R}^\beta_{2\to 1}$. The operators then act within the domains. $\mathcal{E}^\alpha_{2\to 3}\phi(\mathbf{x}) = \int_\Gamma \phi(\mathbf{y})\,\mathcal{K}_\alpha(\mathbf{x} - \mathbf{P}^+\mathbf{y})\,d\mathbf{y}$ is supported in $\mathbf{P}^+\Gamma + \Delta_3(\alpha) \subset \Omega$, and $\mathcal{R}^\alpha_{3\to 2}(\mathbf{1}_\Omega) = 1$ on $\Gamma$. The adjoint relation $\langle u, \mathcal{E}\phi\rangle_{L^2(\Omega)} = \langle \mathcal{R}u, \phi\rangle_{L^2(\Gamma)}$ holds exactly, and similarly for the pair $\Gamma$, $\Sigma$.
\end{definition}

\begin{proposition}[Conservation of Total Mass and Energy Dissipation; formal computation]
Assume \eqref{eq:interface_condition} and homogeneous Neumann conditions on $\partial\Omega$, $\partial\Gamma$, $\partial\Sigma$. Let $(u, v, w)$ be a smooth solution of \eqref{eq:coupled_system}; existence and regularity of solutions are not addressed here. Then:
\begin{enumerate}
    \item \textbf{Conservation of Total Mass:} The total joint mass is strictly invariant for all $t \ge 0$:
    \begin{equation}
    \mathcal{M}_{\text{total}}(t) \coloneqq \int_\Omega u(\mathbf{x}, t)d\mathbf{x} + \int_\Gamma v(\mathbf{y}, t)d\mathbf{y} + \int_\Sigma w(z, t)dz = \mathcal{M}_{\text{total}}(0).
    \end{equation}
    \item \textbf{Global Energy Dissipation:} The joint energy $\mathcal{E}(t) \coloneqq \frac{1}{2}\|u\|_{L^2}^2 + \frac{1}{2}\|v\|_{L^2}^2 + \frac{1}{2}\|w\|_{L^2}^2$ dissipates monotonically:
    \begin{equation}
    \frac{d\mathcal{E}}{dt} = -\|\nabla u\|_{L^2(\Omega)}^2 - \|\nabla_\Gamma v\|_{L^2(\Gamma)}^2 - \|\partial_z w\|_{L^2(\Sigma)}^2 - \lambda \|\mathcal{R}_{3 \to 2}^\alpha u - v\|_{L^2(\Gamma)}^2 - \kappa \|\mathcal{R}_{2 \to 1}^\beta v - w\|_{L^2(\Sigma)}^2 \le 0.
    \end{equation}
\end{enumerate}
\end{proposition}

\begin{proof}
The Laplacian terms integrate to zero by the Neumann conditions. The transmission fluxes between bulk and interface satisfy:
\begin{equation}
\int_\Omega \mathcal{E}_{2 \to 3}^\alpha \big(\lambda (\mathcal{R}_{3 \to 2}^\alpha u - v)\big) d\mathbf{x} = \int_\Gamma \lambda (\mathcal{R}_{3 \to 2}^\alpha u - v)\, \mathcal{R}_{3 \to 2}^\alpha(\mathbf{1}_\Omega)\, d\mathbf{y} = \int_\Gamma \lambda (\mathcal{R}_{3 \to 2}^\alpha u - v) d\mathbf{y},
\end{equation}
since $\mathcal{R}_{3 \to 2}^\alpha(\mathbf{1}_\Omega) = 1$ on $\Gamma$ by the normalization of the kernel and \eqref{eq:interface_condition}. Without \eqref{eq:interface_condition} the factor $\mathcal{R}_{3 \to 2}^\alpha(\mathbf{1}_\Omega)$ is smaller than $1$ near the parts of $\Gamma$ close to $\partial\Omega$, and mass is not conserved. Hence, the mass loss from $\Omega$ identically matches the gain on $\Gamma$. Similarly, the flux between $\Gamma$ and $\Sigma$ cancels, proving $\frac{d\mathcal{M}_{\text{total}}}{dt} = 0$.

For energy dissipation, multiplying by the state fields and applying the adjoint duality $\langle u, \mathcal{E}_{2 \to 3}^\alpha \phi \rangle_{L^2(\Omega)} = \langle \mathcal{R}_{3 \to 2}^\alpha u, \phi \rangle_{L^2(\Gamma)}$ with $\phi = \lambda(\mathcal{R}_{3 \to 2}^\alpha u - v)$:
\begin{align*}
-\langle u, \mathcal{E}_{2 \to 3}^\alpha \phi \rangle_{L^2(\Omega)} + \langle v, \phi \rangle_{L^2(\Gamma)} &= -\langle \mathcal{R}_{3 \to 2}^\alpha u, \phi \rangle_{L^2(\Gamma)} + \langle v, \phi \rangle_{L^2(\Gamma)} \\
&= -\langle \mathcal{R}_{3 \to 2}^\alpha u - v, \lambda(\mathcal{R}_{3 \to 2}^\alpha u - v) \rangle_{L^2(\Gamma)} = -\lambda \|\mathcal{R}_{3 \to 2}^\alpha u - v\|_{L^2(\Gamma)}^2 \le 0.
\end{align*}
Combining with the negative-definite Dirichlet integrals yields monotonic dissipation $\frac{d\mathcal{E}}{dt} \le 0$. All terms are finite for $u \in L^2(\Omega)$ because $\mathcal{R}^\alpha_{3\to 2}$ is bounded on $L^2$ (Remark~\ref{rem:gamma}(c)).
\end{proof}

% ==============================================================================
\section{Reconstruction from Restriction Data}
% ==============================================================================

Let $\operatorname{Gr}(n, m)$ denote the Grassmannian manifold of $n$-dimensional subspaces in $\mathbb{R}^m$. For a projection matrix $\mathbf{P}_\theta$ with $\operatorname{ran}\mathbf{P}_\theta^+ = \theta$, let $g_\theta(\mathbf{z}) = \mathcal{R}_{m \to n, \theta}^\alpha f(\mathbf{z})$ be the observed data. By Definition~\ref{def:radon_beta}, $g_\theta(\mathbf{z}) = h(\mathbf{P}_\theta^+\mathbf{z})$ with $h(\mathbf{x}) \coloneqq \int \mathcal{K}_\alpha(\mathbf{y}) f(\mathbf{x} + \mathbf{y})\,d\mathbf{y}$. The data are therefore \emph{restrictions} of the smoothed field $h$ to planes through the origin, not integrals of $f$ over $n$-planes as in the $n$-plane (Radon) transform.

\begin{proposition}[Inversion]
Let $f \in \mathcal{S}(\mathbb{R}^m)$ and assume $\widehat{\mathcal{K}}_\alpha(\mathbf{k}) \ne 0$ for all $\mathbf{k}$. For every $\mathbf{x} \in \mathbb{R}^m$ and every $\theta \ni \mathbf{x}$,
\begin{equation}
h(\mathbf{x}) = g_\theta(\mathbf{P}_\theta \mathbf{x}),
\end{equation}
so the family $\{g_\theta\}$ determines $h$ on $\bigcup_\theta \theta = \mathbb{R}^m$. Then $\widehat{h}(\mathbf{k}) = \widehat{\mathcal{K}}_\alpha(-\mathbf{k})\widehat{f}(\mathbf{k})$, and
\begin{equation}
f = \mathcal{F}^{-1}\Big[\widehat{h}(\mathbf{k}) / \widehat{\mathcal{K}}_\alpha(-\mathbf{k})\Big].
\end{equation}
\end{proposition}
\begin{proof}
For $\mathbf{x} \in \theta = \operatorname{ran}\mathbf{P}_\theta^+ = \operatorname{ran}\mathbf{P}_\theta^T$, $\mathbf{P}_\theta^+\mathbf{P}_\theta$ is the orthogonal projection onto $\theta$, so $\mathbf{P}_\theta^+(\mathbf{P}_\theta\mathbf{x}) = \mathbf{x}$. The Fourier identity is the convolution theorem for $h = \check{\mathcal{K}}_\alpha * f$ with $\check{\mathcal{K}}_\alpha(\mathbf{y}) = \mathcal{K}_\alpha(-\mathbf{y})$.
\end{proof}

\begin{remark}[Filtered backprojection does not apply]
A filtered backprojection with filter $|\boldsymbol{\xi}|^{m-n}$ inverts the $n$-plane transform, because of the Fourier slice theorem. Here the Fourier transform of the data is instead an integral of $\widehat{h}$ over the fibres $\mathbf{P}_\theta^T\boldsymbol{\xi} + \ker\mathbf{P}_\theta$ (Proposition~\ref{prop:fourier_rep} with a $\delta$ kernel), so the slice theorem does not hold. Numerically, for $m = 2$, $n = 1$, a Gaussian $f$ centred at $(1.5, 0)$ and $\alpha = 1$, the filtered backprojection of the restriction data divided by $h$ takes the values $11.28$, $10.01$, $-0.90$ and $14.73$ at four test points. Applied to genuine line integrals, the same procedure gives the constant $(2\pi)^2$. A tomographic model based on the Beta-kernel would have to integrate over the fibres of $\mathbf{P}_\theta$. That is a different operator from $\mathcal{R}^\alpha_{m\to n}$.
\end{remark}

\begin{remark}[Ill-posedness]
Division by $\widehat{\mathcal{K}}_\alpha$ is a deconvolution. Because $\widehat{\mathcal{K}}_\alpha$ decays at high frequency, the division \emph{amplifies} high-frequency noise, and at any zero of $\widehat{\mathcal{K}}_\alpha$ it is undefined. The reconstruction is therefore ill-posed and requires regularization, for instance Tikhonov filtering $\overline{\widehat{\mathcal{K}}_\alpha}/(|\widehat{\mathcal{K}}_\alpha|^2 + \epsilon)$. By Remark~\ref{rem:gamma}, $|\widehat{\mathcal{K}}_\alpha|$ decays like $|\mathbf{k}|^{-1}$ in facet-normal directions, so the deconvolution is mildly ill-posed of order one. The kernel therefore does not regularize the inversion or suppress Gibbs ringing; the effect is the opposite. The Beta-kernel smooths the \emph{forward} projections; that smoothing is exactly what makes the inverse problem harder.
\end{remark}

% ==============================================================================
\section{A Barycentric Contraction Map}
% ==============================================================================

In high-dimensional data analysis, multi-way relational structures are encoded as abstract simplicial complexes $K = \{\sigma_1, \dots, \sigma_N\}$ with barycentric coordinates $\mathbf{x} = \sum_{j=1}^m c_j \mathbf{v}_j$, where $\sum c_j = 1, c_j \ge 0$.

\begin{definition}[Induced Barycentric Contraction Map]
\label{def:induced_barycentric_map}
Let $\sigma$ be the $(m-1)$-simplex with vertex set $\{1,\dots,m\}$ and let $\pi : \{1,\dots,m\} \twoheadrightarrow \{1,\dots,n\}$ be a surjective grouping of the vertices, chosen by the user, with fibers $G_k \coloneqq \pi^{-1}(k)$, $k=1,\dots,n$. On the barycentric coordinate vector $\mathbf{c} = (c_1,\dots,c_m) \in \Delta^{m-1}$ of a point $\mathbf{x} = \sum_j c_j \mathbf{v}_j \in \sigma$ define $\Phi_\alpha : \Delta^{m-1} \to \Delta^{n-1}$ by
\begin{equation}
\label{eq:induced_map}
\Phi_\alpha(\mathbf{c})_k \coloneqq \frac{\alpha \sum_{j \in G_k} c_j + |G_k|}{\alpha + m}, \qquad k = 1, \dots, n,
\end{equation}
the mean of the Dirichlet law $\mathrm{Dir}(\alpha c_1 + 1, \dots, \alpha c_m + 1)$ pushed forward along $\pi$. The Dirichlet density is the discrete-parameter analogue of the Beta-kernel. The map $\Phi_\alpha$ is a separate construction: it is not derived from the transform $\mathcal{R}_{m\to n}^\alpha$ of Definition~\ref{def:radon_beta}, and a rank-$n$ projection does not in general induce a vertex grouping. Note $\sum_k \Phi_\alpha(\mathbf{c})_k = \frac{\alpha \sum_j c_j + m}{\alpha+m} = 1$, so $\Phi_\alpha(\mathbf{c}) \in \Delta^{n-1}$ is well-defined.
\end{definition}

\begin{theorem}[Barycentric Lipschitz Bound with Sharp Constant]
\label{thm:barycentric_ratio}
Let $\sigma \in K$ be an $(m-1)$-simplex and $\Phi_\alpha : \Delta^{m-1} \to \Delta^{n-1}$ the induced map of Definition~\ref{def:induced_barycentric_map}. For any $\mathbf{c}_1, \mathbf{c}_2 \in \Delta^{m-1}$,
\begin{equation}
\label{eq:sharp_contraction}
\operatorname{dist}_\Delta\big(\Phi_\alpha(\mathbf{c}_1), \Phi_\alpha(\mathbf{c}_2)\big) \;\le\; \frac{\alpha}{\alpha+m}\,\operatorname{dist}_\Delta(\mathbf{c}_1, \mathbf{c}_2),
\end{equation}
where $\operatorname{dist}_\Delta(\mathbf{c},\mathbf{c}') \coloneqq \tfrac{1}{2}\|\mathbf{c}-\mathbf{c}'\|_1$ is the barycentric Wasserstein-1 metric on the equilateral vertex embedding. The constant $\alpha/(\alpha+m)$ is sharp: equality holds whenever $\mathbf{c}_1 - \mathbf{c}_2$ does not change sign within any fiber $G_k$ (in particular whenever $n = m$, i.e.\ $\pi$ is the identity). Consequently
\begin{equation}
\frac{\operatorname{dist}_\Delta(\Phi_\alpha \mathbf{c}_1, \Phi_\alpha \mathbf{c}_2)}{\operatorname{dist}_\Delta(\mathbf{c}_1, \mathbf{c}_2)} \;\le\; 1 - \frac{m}{\alpha} + \mathcal{O}(\alpha^{-2}) \qquad (\alpha \to \infty),
\end{equation}
with the ratio identically $0$ whenever $\pi$ merges vertices on which $\mathbf{c}_1,\mathbf{c}_2$ agree oppositely within every fiber (e.g.\ $m=4 \to n=2$, $G_1=\{1,2\}$, $G_2=\{3,4\}$, $\mathbf{c}_1=(0.5,0,0.5,0)$, $\mathbf{c}_2=(0,0.5,0,0.5)$) --- so the ratio does \emph{not} converge to $1$ uniformly over $\Delta^{m-1}\times\Delta^{m-1}$ for $n<m$, and there is no lower (bi-Lipschitz) bound. The map $\Phi_\alpha$ is affine, so \eqref{eq:sharp_contraction} is an upper Lipschitz bound obtained from the triangle inequality. It is not a two-sided distortion estimate.
\end{theorem}

\begin{proof}
Fix $k \in \{1,\dots,n\}$. From \eqref{eq:induced_map},
\begin{equation}
\Phi_\alpha(\mathbf{c}_1)_k - \Phi_\alpha(\mathbf{c}_2)_k = \frac{\alpha}{\alpha+m} \sum_{j \in G_k} (c_{1,j} - c_{2,j}).
\end{equation}
By the triangle inequality, $\left|\sum_{j \in G_k}(c_{1,j}-c_{2,j})\right| \le \sum_{j\in G_k}|c_{1,j}-c_{2,j}|$, with equality iff $c_{1,j}-c_{2,j}$ has constant sign on $G_k$. Summing over $k=1,\dots,n$ and using that $\{G_k\}$ partitions $\{1,\dots,m\}$:
\begin{equation}
\begin{aligned}
\|\Phi_\alpha(\mathbf{c}_1) - \Phi_\alpha(\mathbf{c}_2)\|_1 &= \frac{\alpha}{\alpha+m}\sum_{k=1}^n \left|\sum_{j\in G_k}(c_{1,j}-c_{2,j})\right| \\ &\le \frac{\alpha}{\alpha+m}\sum_{k=1}^n \sum_{j\in G_k}|c_{1,j}-c_{2,j}| = \frac{\alpha}{\alpha+m}\|\mathbf{c}_1-\mathbf{c}_2\|_1.
\end{aligned}
\end{equation}
Dividing by $2$ gives \eqref{eq:sharp_contraction}. Equality throughout requires equality in every fiber's triangle inequality, i.e.\ no sign cancellation within any $G_k$; when $n=m$ every fiber is a singleton and this holds automatically, giving the exact identity $\operatorname{dist}_\Delta(\Phi_\alpha\mathbf{c}_1,\Phi_\alpha\mathbf{c}_2) = \frac{\alpha}{\alpha+m}\operatorname{dist}_\Delta(\mathbf{c}_1,\mathbf{c}_2)$. The stated counterexample exhibits full cancellation ($\sum_{j\in G_k}(c_{1,j}-c_{2,j})=0$ for every $k$), forcing $\Phi_\alpha(\mathbf{c}_1)=\Phi_\alpha(\mathbf{c}_2)$ identically in $\alpha$, which shows the ratio bound \eqref{eq:sharp_contraction} cannot be improved to an asymptotic equality valid uniformly over $\Delta^{m-1}\times\Delta^{m-1}$. Finally, $\frac{\alpha}{\alpha+m} = 1 - \frac{m}{\alpha+m} = 1-\frac{m}{\alpha}+\mathcal{O}(\alpha^{-2})$ as $\alpha\to\infty$ by direct Taylor expansion.
\end{proof}

% ==============================================================================
\section{Matrix-Valued Beta-Kernels}
% ==============================================================================

Let $\operatorname{Sym}_m^+(\mathbb{R})$ denote the open convex cone of $m \times m$ real symmetric positive-definite matrices. This section recalls classical facts, which we do not extend here, for instance to Grassmannians.

\begin{definition}[Siegel--Wishart Matrix Beta Operator]
Extending multivariate hypergeometric distributions and Wishart random matrices \cite{muirhead2009}, for scalar shape parameters $a, b \in \mathbb{R}$ with $a, b > \frac{m-1}{2}$, the matrix Beta operator acting on $F: \operatorname{Sym}_m^+(\mathbb{R}) \to \mathbb{R}$ is defined by
\begin{equation}
\mathcal{G}_{a, b} F(\mathbf{X}) \coloneqq \frac{1}{\mathrm{B}_m(a, b)} \int_{\mathbf{0} < \mathbf{Y} < \mathbf{I}_m} (\det \mathbf{Y})^{a - \frac{m+1}{2}} (\det(\mathbf{I}_m - \mathbf{Y}))^{b - \frac{m+1}{2}} F\left(\mathbf{X}^{1/2} \mathbf{Y} \mathbf{X}^{1/2}\right) d\mathbf{Y},
\end{equation}
where the multivariate Siegel Beta normalization is given by
\begin{equation}
\mathrm{B}_m(a, b) \coloneqq \frac{\Gamma_m(a)\Gamma_m(b)}{\Gamma_m(a+b)}, \quad \Gamma_m(a) \coloneqq \pi^{m(m-1)/4}\prod_{j=1}^m \Gamma\left(a - \frac{j-1}{2}\right).
\end{equation}
\end{definition}

\begin{theorem}[Invariance and Spherical Harmonic Eigenfunctions; classical, see {\cite[Ch.~7]{muirhead2009}}]
The operator $\mathcal{G}_{a, b}$ is invariant under the congruence action of the orthogonal group $\mathrm{O}(m)$ ($\mathbf{X} \mapsto \mathbf{U}\mathbf{X}\mathbf{U}^T$). Furthermore, for any zonal spherical polynomial $Z_\lambda(\mathbf{X})$ associated with partition $\lambda$:
\begin{equation}
\mathcal{G}_{a, b} Z_\lambda(\mathbf{X}) = \frac{(a)_\lambda}{(a+b)_\lambda} Z_\lambda(\mathbf{X}),
\end{equation}
where $(a)_\lambda \coloneqq \prod_{j=1}^m \frac{\Gamma\left(a - \frac{j-1}{2} + \lambda_j\right)}{\Gamma\left(a - \frac{j-1}{2}\right)}$ denotes the multivariate generalized Pochhammer symbol.
\end{theorem}

% ==============================================================================
\section{Conclusion}
% ==============================================================================

We have studied dimension-changing non-local operators $\mathcal{T}_{m \to n}^\alpha$. We proved a Sobolev trace estimate whose smoothing gain is the Fourier decay rate of the kernel, which equals $1$ for the Beta-kernel, and gave the dual extension estimate. We formulated a coupled 3D--2D--1D interface system, conservative under an explicit interface condition (formal computation). We showed that the transform yields restrictions of a smoothed field, so reconstruction is a deconvolution, which is ill-posed. We proved a Lipschitz bound with sharp constant for a barycentric Dirichlet-mean map, and recalled the Siegel--Wishart matrix Beta operators. Determining the exact directional Fourier decay of the Beta-kernel beyond the isotropic exponent is left open. Future work will investigate numerical finite-element schemes for inter-dimensional multiscale PDEs.

% Shared declarations block, \input by every chapter before its bibliography.
% Keep in sync with the "Declarations" section of master_book_unified_quantum_gravity.tex.
\section*{Declarations}

\noindent\textbf{Funding.} This research was financed in part by the Coordena\c{c}\~ao de Aperfei\c{c}oamento de Pessoal de N\'ivel Superior -- Brasil (CAPES) -- Finance Code 001.

\medskip\noindent\textbf{Competing interests.} The author declares no competing interests.

\medskip\noindent\textbf{Use of generative AI tools.} Large language model assistants (Anthropic Claude; Google Gemini, used through the Antigravity agent environment) were used extensively in preparing this work: drafting and editing text and \LaTeX{} source; proposing, expanding and revising derivations and proofs under the author's direction; writing the Python verification scripts and the \textsc{Lean 4} files; searching and checking the literature and references; and adversarial auditing of proofs, numerical claims and code. These audits found errors, some of them originally introduced with AI assistance. The errors and their corrections are recorded in the repository file \texttt{WORKPLAN.md}. AI tools are not authors. The author chose the research questions and the overall framework, reviewed the material, and is responsible for the content, including any remaining errors.

\medskip\noindent\textbf{Verification status.} Results labelled Theorem, Proposition or Lemma are proved in the text or cited as classical; statements labelled Conjecture are not proved. The numerical scripts compare formulas of the text with independent computations and are checks, not proofs. The \textsc{Lean 4} modules in \texttt{formal\_proofs\_book/Book} are a naming skeleton and do not verify the mathematics of this chapter.

\medskip\noindent\textbf{Code and data availability.} Sources, numerical scripts and the audit log are available at
\begin{center}\url{https://github.com/reinaldomsilvafilho-netizen/quantum-gravity-lean4}.\end{center}
The monograph is archived on Zenodo, \href{https://doi.org/10.5281/zenodo.22290043}{doi:10.5281/zenodo.22290043}, under the CC~BY~4.0 license.


\begin{thebibliography}{99}
\bibitem{silvafilho2026simplex} R.~M. Silva-Filho, \emph{Analytic Continuation of Pascal's Simplex: Continuous Multinomial Integrals, Polytope Boundary Recurrences, and Simplicial Fractional Calculus}, Chapter~3 of \emph{Geometry, Tensors, and Quantum Gravity}, Universidade Federal de Lavras, 2026.
\bibitem{silvafilho2026waves} R.~M. Silva-Filho, \emph{Nonlinear Simplicial Waves and Anomalous Porous Transport Induced by Beta-Kernel Fractional Laplacians}, Chapter~4 of \emph{Geometry, Tensors, and Quantum Gravity}, Universidade Federal de Lavras, 2026.
\bibitem{helgason2011} S. Helgason, \emph{Integral Geometry and Radon Transforms}, Springer, New York, 2011, \href{https://doi.org/10.1007/978-1-4419-6055-9}{doi:10.1007/978-1-4419-6055-9}.
\bibitem{samko1993} S. G. Samko, A. A. Kilbas, and O. I. Marichev, \emph{Fractional Integrals and Derivatives: Theory and Applications}, Gordon and Breach, Yverdon, 1993, ISBN 978-2-88124-864-1.
\bibitem{adams2003} R. A. Adams and J. J. F. Fournier, \emph{Sobolev Spaces}, 2nd ed., Pure and Applied Mathematics, vol. 140, Academic Press, New York, 2003, ISBN 978-0-12-044143-3.
\bibitem{muirhead2009} R. J. Muirhead, \emph{Aspects of Multivariate Statistical Theory}, John Wiley \& Sons, Hoboken, 2009, \href{https://doi.org/10.1002/9780470316559}{doi:10.1002/9780470316559}.
\bibitem{evans2010} L. C. Evans, \emph{Partial Differential Equations}, 2nd ed., Graduate Studies in Mathematics, vol. 19, American Mathematical Society, Providence, 2010, \href{https://doi.org/10.1090/gsm/019}{doi:10.1090/gsm/019}.
\end{thebibliography}

\end{document}
